% ============================================================================
%  CAPITULO IV: CONCLUSIONES
% ============================================================================

\section{Conclusiones}

El lenguaje de control de datos (DCL) fue el más difícil de ejecutar, porque:

\begin{enumerate}
  \item La administración de permisos es más amplia y detallada que el resto de lenguajes: hay que crear el rol (\texttt{CREATE ROLE}), asignar \texttt{GRANT}s, revocarlos (\texttt{REVOKE}) y verificar el acceso real desde una conexión externa, viendo tanto lo permitido como lo denegado.
  \item El manejo de los permisos por defecto de \texttt{PUBLIC} y la confirmación de las denegaciones añaden pasos extra (ej. \texttt{REVOKE} sobre \texttt{PUBLIC} para asegurar la denegación).
\end{enumerate}

Respecto a los demás lenguajes:

\begin{enumerate}
  \item DDL: el más sencillo de formular; la parte más compleja fue la semántica de \texttt{ALTER}/\texttt{DROP} sobre objetos que ya tienen dependencias.
  \item DML: directo; los \texttt{UPDATE}/\texttt{DELETE} con subconsultas en \texttt{WHERE} fueron sencillos de construir.
  \item TCL: claro; el uso de \texttt{SAVEPOINT} permite revertir una parte de la transacción sin perder la \texttt{COMMIT}.
\end{enumerate}
